\section{Conclusion}

This article extends the previously verified foundations for recursive
lists, discrete integrals, modulo arithmetic, and cycles to define and
reason about Cycle Integrals. Starting from a finite non-empty list, the
construction treats the list as a repeating cycle and describes the
accumulated value at any non-negative index using the cycle sum, modular
position, and initial value.

We defined two equivalent presentations of Cycle Integral:
\textbf{CycleIntegral}, a recursive accumulation over a memory-backed cycle
(Subsection~3.1), and \textbf{ModCycleIntegral}, a closed-form definition
using division and modulo (Subsection~3.2).

For both presentations, we verified the sum property (integral equals
cumulative cycle sum) and the step property (difference between consecutive
values equals the corresponding cycle element). We also proved equivalence
of the recursive and modulo definitions, that the integral is strictly
increasing under positive base values, and that it stays positive
everywhere given a non-negative start.

Beyond these core definitions, the article verifies several reusable
cycle-integral laws. Within a fixed cycle integral: residues are periodic
when the cycle sum is zero modulo the chosen modulus, with the persistent
non-zero and persistent zero cases following as immediate corollaries;
adjacent gaps telescope into integral differences; full-cycle shifts advance
the integral by the cycle sum; and cycle residue classification is correct,
exclusive, and exhaustive. Building new cycle integrals from existing ones:
repeating a backing cycle preserves the represented integral stream;
rotating a gap cycle with the corresponding head adjustment shifts the
represented integral by one position; survivor scans retain exactly the
non-multiples needed for filtering; and merging the two gaps around a
removed multiple reconstructs exactly the cycle integral that a fresh
construction from the survivor list would produce, with no multiples of
the filter value remaining anywhere in the result. The open extension
points are the index-shift laws in Sections 6.2 and 6.3.

The main established properties are:

\begin{equation*}
\begin{aligned}
&\forall \ L \in \mathbb{L},\quad \forall \ init \in \mathbb{N}_0,\quad \forall \ i \in \mathbb{N}_0 \\
&L = [v_0, v_1, \dots, v_{n-1}], \quad n = |L|,\quad n > 0 \\
&T = \sum_{j=0}^{n-1} v_j
\end{aligned}
\end{equation*}

\begin{equation*}
\operatorname{CycleIntegral}(L, init)_i
= \left(i \operatorname{div} n\right) \cdot T
 + \operatorname{CycleIntegral}(L, init)_{i \operatorname{mod} n}
\quad \text{[Modulo Cycle Integral]}
\end{equation*}

\begin{equation*}
\operatorname{CycleIntegral}(L, init)_i
= \operatorname{ModCycleIntegral}(L, init)_i
\quad \text{[Definition Equivalence]}
\end{equation*}

\begin{equation*}
\operatorname{CycleIntegral}(L, init)_{i+1}
- \operatorname{CycleIntegral}(L, init)_i
= \operatorname{Cycle}(L)_{i+1}
\quad \text{[Step Property]}
\end{equation*}

\begin{equation*}
\begin{aligned}
\operatorname{CycleIntegral}(L, init)_{i+1}
&- \operatorname{CycleIntegral}(L, init)_i \\
&= \operatorname{CycleIntegral}(L, init)_{i+n+1}
 - \operatorname{CycleIntegral}(L, init)_{i+n} \\
&\quad \text{[Same Difference After Full Cycle]}
\end{aligned}
\end{equation*}

\begin{equation*}
\operatorname{CycleIntegral}(L, init)_i
= \operatorname{sum}([init] \mathbin{\texttt{++}}
 [\operatorname{Cycle}(L)_0, \ldots, \operatorname{Cycle}(L)_i])
\quad \text{[Sum of Mod Values as List]}
\end{equation*}

\begin{equation*}
\begin{aligned}
init \geq 0 \land (\forall x \in L,\ x > 0) \land b > a
&\implies
\operatorname{CycleIntegral}(L, init)_b > \operatorname{CycleIntegral}(L, init)_a \\
&\quad \text{[Strictly Increasing]}
\end{aligned}
\end{equation*}

\begin{equation*}
init \geq 0 \land (\forall x \in L,\ x > 0)
\implies
\operatorname{CycleIntegral}(L, init)_i > 0
\quad \text{[Positivity]}
\end{equation*}

\begin{equation*}
ci(pos + \operatorname{period}(ci))
= ci(pos) + \operatorname{periodSum}(ci)
\quad \text{[Cycle-Period Shift]}
\end{equation*}

\begin{equation*}
\begin{aligned}
\operatorname{mod}(\operatorname{periodSum}(ci), m) = 0
&\implies
\operatorname{mod}(ci(pos), m) = \operatorname{mod}(ci(pos \bmod n), m) \\
&\quad \text{[General Residue Periodicity]}
\end{aligned}
\end{equation*}

\begin{equation*}
\begin{aligned}
\operatorname{mod}(&\operatorname{periodSum}(ci), m) = 0 \land \big(\forall\, k \in [0,n),\ \operatorname{mod}(ci(k), m) \neq 0\big) \\
&\implies
\forall\, i,\ \operatorname{mod}(ci(i), m) \neq 0
\quad \text{[Persistent Non-Zero Residue]}
\end{aligned}
\end{equation*}

\begin{equation*}
\begin{aligned}
\operatorname{mod}(&\operatorname{periodSum}(ci), m) = 0 \land \big(\forall\, k \in [0,n),\ \operatorname{mod}(ci(k), m) = 0\big) \\
&\implies
\forall\, i,\ \operatorname{mod}(ci(i), m) = 0
\quad \text{[Persistent Zero Residue]}
\end{aligned}
\end{equation*}

\begin{equation*}
\operatorname{CycleIntegral}(G,h)_{k+1}
- \operatorname{CycleIntegral}(G,h)_{k-1}
= G_{k-1} + G_k
\quad \text{[Two-Gap Telescoping]}
\end{equation*}

\begin{equation*}
\operatorname{classify}(I_i,m)
\in \{\text{zero},\text{nonzero}\}
\quad \text{[Residue Classification]}
\end{equation*}

\begin{equation*}
\operatorname{CycleIntegral}(L^{\langle x\rangle}, init)_i
= \operatorname{CycleIntegral}(L, init)_i
\quad \text{[Repeated-Cycle Invariance]}
\end{equation*}

\begin{equation*}
\operatorname{rotateAt}(G,1)\text{ with head }h+G_0
\implies
I'_{i}=I_{i+1}
\quad \text{[Rotation Shift]}
\end{equation*}

\begin{equation*}
\operatorname{survivors}(I,m)
= \{ I_i \mid I_i \not\equiv 0 \pmod m \}
\quad \text{[Survivor Exactness]}
\end{equation*}

\begin{equation*}
\begin{aligned}
&\big(\forall\, q \in [start,pos),\ \operatorname{mod}(ci(q),f)=0\big) \land \operatorname{mod}(ci(pos),f)\neq 0 \\
&\quad \implies
\operatorname{head}(S) = ci(pos)
\quad \text{[Survivor Structure]}
\end{aligned}
\end{equation*}

\begin{equation*}
ci'(m) = ci(m + 1)
\quad \text{[Merge Shift Law]}
\end{equation*}

\begin{equation*}
\operatorname{mod}(ci(m), f) = 0 \implies ci'(m) = ci(m + 1)
\quad \text{[Removing a Multiple]}
\end{equation*}

\begin{equation*}
\begin{aligned}
ci''\text{'s initial value} = S_0 \land \operatorname{cycle}''(i) = S_{i+1} - S_i
&\implies ci''(k) = S_{k+1} \\
&\quad \text{[Direct Construction from Survivors]}
\end{aligned}
\end{equation*}

\begin{equation*}
\operatorname{mod}(ci(0), f) \neq 0 \implies \operatorname{mod}(ci''(k), f) \neq 0
\quad \text{[Filtered Result Has No Multiples]}
\end{equation*}

The verified definitions provide a reusable foundation for reasoning about
infinite periodic accumulations using finite list structures and
machine-checked Scala code.

\section{Future Work}

The nearest continuation is to close the remaining index-shift properties in
Section 6. Those statements already have mathematical derivations in the
article, but they still sit at the boundary between the verified
cycle-integral library and the stronger shift reasoning needed for later
sieve arguments.

After that, the same finite-period machinery can support more specialized
prime-sieve properties: detecting accepted residues, tracking gap evolution
across filters, and relating local survivor windows to complete-cycle
structure. More distant extensions, such as multi-dimensional cycles or
integration over richer algebraic structures, would require new definitions
rather than a direct continuation of the present proof.
